% Requires graphicx. Set \EDAfigpath if compiling outside the repository root.
\providecommand{\EDAfigpath}{paper/figures}
\subsection{Exploratory data analysis}
\label{sec:eda}

Exploratory analysis covered the 65,287 eligible training households. It characterized recorded trip counts, examined household attributes relevant to trip generation, and assessed the sample support for the proposed cross-classification cells.

Recorded weekday trips varied substantially across households (Table~\ref{tab:eda-summary}). The mean was 5.64 trips, the median was four, and the standard deviation was 4.55. The middle half of households reported between two and eight trips; the 95th percentile was 14, and the maximum was 55. The distribution concentrated at low counts and had a long upper tail (Figure~\ref{fig:eda-trips}). It included 7,274 households (11.14\%) with no recorded trips. Retaining these households preserved the observed range of weekday travel activity. The tail and the zero-trip mass warranted attention to residual shape, large prediction errors, and negative OLS predictions.

\begin{figure}[!htbp]
\centering
\includegraphics[width=0.9\linewidth]{\EDAfigpath/eda_trip_distribution.pdf}
\caption{Distribution of recorded weekday trips among eligible training households. Counts of 25 or more are grouped for display; the descriptive statistics use the original counts.}
\label{fig:eda-trips}
\end{figure}

\begin{table}[htbp]
\centering
\caption{Selected characteristics of eligible training households ($n=65{,}287$).}
\label{tab:eda-summary}
\begin{tabular}{lrrrr}
\hline
Characteristic & Mean & SD & Median & 95th percentile \\
\hline
Recorded weekday trips & 5.64 & 4.55 & 4 & 14 \\
Household persons & 2.74 & 1.42 & 2 & 5 \\
Vehicles & 1.38 & 0.93 & 1 & 3 \\
Licensed drivers & 1.73 & 0.98 & 2 & 3 \\
Full-time workers & 1.11 & 0.93 & 1 & 3 \\
Students & 0.63 & 0.95 & 0 & 3 \\
\hline
\end{tabular}
\end{table}

Mean trip rates increased across the four characteristics in the starting regression (Figure~\ref{fig:eda-rates}). One-person households averaged 2.04 trips, compared with 10.16 among households with five or more persons. Households with no vehicle averaged 2.35 trips, rising to 9.04 for those with three or more vehicles. The corresponding means were 3.04 and 7.53 trips for households with zero and two full-time workers, and 3.99 and 9.25 for households with zero and two students. These marginal patterns accord with the roles of household size, vehicle availability, employment, and education-related activity in trip generation. Households with more workers or students also tend to have more residents, so the regression estimates their adjusted associations. The student association requires particular care because the student count includes children aged 6--10, whose trips fall outside the recorded outcome.

\begin{figure}[!htbp]
\centering
\includegraphics[width=\linewidth]{\EDAfigpath/eda_predictor_rates.pdf}
\caption{Unadjusted mean weekday trip rates across the four starting regression predictors. A plus sign denotes a pooled upper category. The common vertical scale permits comparison of the profiles; the lines connect category means and are not fitted regression curves.}
\label{fig:eda-rates}
\end{figure}

Licensed drivers added useful variation within household-size groups (Figure~\ref{fig:eda-cells}). Among four-person households, mean trips increased from 4.80 with no licensed driver to 9.73 with three or more. This pattern motivated the joint size--driver classification and the inclusion of driver count among the regression alternatives. Grouping size as 1, 2, 3, 4, and 5 or more and drivers as 0, 1, 2, and 3 or more yielded 17 feasible cells. All contained observations, with counts from 96 to 13,327, and exceeded the minimum support threshold of 30. Three combinations were structurally impossible. The variation in cell counts remained relevant to the precision of the estimated rates.

\begin{figure}[!htbp]
\centering
\includegraphics[width=0.95\linewidth]{\EDAfigpath/eda_size_driver_cells.pdf}
\caption{Training mean weekday trips and household counts for size--driver cells. Each observed cell displays its mean rate and sample size $n$. Blank grey cells are structurally impossible. Upper categories pool households with five or more persons or three or more drivers.}
\label{fig:eda-cells}
\end{figure}

The group-rate profiles describe how trips vary with each household characteristic; the correlation heatmap (Figure~\ref{fig:eda-correlations}) shows how the characteristics vary together. Pearson coefficients summarize linear association on the original count scale used by OLS. Vehicles and licensed drivers were strongly correlated ($r=0.713$), as were household size and students ($r=0.689$). These overlaps motivated examination of adjusted coefficients and variance-inflation factors in the fitted models. Total workers was considered as an alternative to its full-time and part-time components, whose sum is exactly the total. Correlations with trips remain marginal descriptions. Their sensitivity to extreme trip counts reinforces the use of group profiles and influence diagnostics alongside the heatmap \cite{deWinter2016}.

\begin{figure}[!htbp]
\centering
\includegraphics[width=0.92\linewidth]{\EDAfigpath/eda_correlation_heatmap.pdf}
\caption{Pearson correlations among recorded trips and selected household counts in the eligible training sample ($n=65{,}287$). The full symmetric matrix is shown; total workers is omitted because it is the sum of full-time and part-time workers.}
\label{fig:eda-correlations}
\end{figure}

Trip rates also differed by dwelling, region, and survey weekday. Mean trips were 6.56 for houses, 3.64 for apartments, and 6.02 for townhouses. Regional means ranged from 4.93 in Toronto to 6.70 in York, while weekday means ranged from 5.43 on Monday to 5.87 on Friday. Differences in household composition could contribute to these contrasts; the regression search therefore assessed dwelling, region, and weekday as categorical effects.

The expansion-weighted training mean was 5.6385 trips per household, close to the unweighted mean of 5.6424. Overall agreement leaves open differences within groups and zones, which the weighted evaluation examines separately. Upper categories such as 5+ persons and 3+ drivers pooled observations for descriptive and cross-classification summaries; regression used the original eligible counts.

\subsubsection{Methodological decisions informed by the exploratory analysis}

The starting regression retains the four prespecified household characteristics, while the alternative specifications address the observed overlap among household attributes and the variation across residential settings. The skewed trip distribution calls for both MAE and RMSE, along with residual and influential-observation checks. The 7,194 households with five or more persons and 10,443 with three or more licensed drivers supported pooling the upper categories in the primary cross-classification. Every feasible cell in its 5+/3+ grid had at least 96 training households. The 4+/3+ and 5+/4+ grids provide sensitivity checks on those boundaries.
